% ---------------------------------------------------------------------------
% ---------------------------------------------------------------------------
% Modelo LaTex para preparação do documento final de Monografia TCC
% O modelo está em conformidade com ABNT NBR
% Unidade de Ensino Superior Dom Bosco - UNDB
% Autor: Prof. Me. Allan Kássio Beckman Soares da Cruz
% Curso de Sistemas da Informação
% ---------------------------------------------------------------------------
% ---------------------------------------------------------------------------

\documentclass[
	% -- opções da classe memoir --
	12pt,						% tamanho da fonte
	openright,					% capítulos começam em pág ímpar (insere página vazia caso preciso)
	oneside,					% para impressão em verso e anverso. Oposto a oneside
	a4paper,					% tamanho do papel. 
	% -- opções da classe abntex2 --
	chapter=TITLE,				% títulos de capítulos convertidos em letras maiúsculas
	%section=TITLE,				% títulos de seções convertidos em letras maiúsculas
	%subsection=TITLE,			% títulos de subseções convertidos em letras maiúsculas
	%subsubsection=TITLE,		% títulos de subsubseções convertidos em letras maiúsculas
    sumario=abnt-6027-2012,
	% -- opções do pacote babel --
	english,					% idioma adicional para hifenização
	%french,					% idioma adicional para hifenização
	%spanish,					% idioma adicional para hifenização
	brazil						% o último idioma é o principal do documento
	]{abntex2}

% ---------------------
% Pacotes OBRIGATÓRIOS
% ---------------------
%\usepackage{lmodern}				% Usa a fonte Latin Modern	
\usepackage{uarial}					% usa fonte arial
\renewcommand*\familydefault{\sfdefault}
\usepackage[T1]{fontenc}			% Selecao de codigos de fonte.
\usepackage[utf8]{inputenc}			% Codificacao do documento (conversão automática dos acentos)
\usepackage{lastpage}				% Usado pela Ficha catalográfica
\usepackage{indentfirst}			% Indenta o primeiro parágrafo de cada seção.
\usepackage{color}					% Controle das cores
\usepackage{graphicx,graphicx}		% Inclusão de gráficos
\usepackage{epsfig,subfig}			% Inclusão de figuras
\usepackage{microtype} 				% Melhorias de justificação
% ---------------------
		
% ---------------------
% Pacotes ADICIONAIS
% ---------------------
\usepackage{lipsum}						% Geração de dummy text
\usepackage{amsmath,amssymb,mathrsfs}	% Comandos matemáticos avançados 
\usepackage{setspace}  					% Para permitir espaçamento simples, 1 1/2 e duplo
\usepackage{verbatim}					% Para poder usar o ambiente "comment"
\usepackage{tabularx} 					% Para poder ter tabelas com colunas de largura auto-ajustável
\usepackage{afterpage} 					% Para executar um comando depois do fim da página corrente
\usepackage{url} 						% Para formatar URLs (endereços da Web)
% ---------------------

\renewcommand{\ABNTEXchapterfontsize}{\bfseries{\normalsize}}
\renewcommand{\ABNTEXpartfontsize}{\bfseries{\normalsize}}
\renewcommand{\ABNTEXsectionfontsize}{\bfseries{\normalsize}}
\renewcommand{\ABNTEXsubsectionfontsize}{\normalsize}
\renewcommand{\ABNTEXsubsubsectionfontsize}{\normalsize}
\renewcommand{\ABNTEXsubsubsubsectionfontsize}{\normalsize}
\setlength\afterchapskip{\lineskip}




% ---------------------
% Pacotes de CITAÇÕES
% ---------------------
\usepackage[brazilian,hyperpageref]{backref}	% Paginas com as citações na bibl
\usepackage[alf]{abntex2cite}					% Citações padrão ABNT (alfa)
%\usepackage[num]{abntex2cite}					% Citações padrão ABNT (numericas)
% ---------------------

% Configurações de CITAÇÕES para abntex2
\include{extras/conf_citacoes}

% Inclusão de dados para CAPA e FOLHA DE ROSTO (título, autor, orientador, etc.)
\include{extras/dados}

% Inclui Configurações de aparência do PDF Final
\include{extras/conf_pdf}

% O tamanho da identação do parágrafo é dado por:
\setlength{\parindent}{1.3cm}

% Controle do espaçamento entre um parágrafo e outro:
\setlength{\parskip}{0.2cm}  % tente também \onelineskip

% ---------------------
% Compila o indice
% ---------------------
\makeindex
% ---------------------

%%%%%%%%%%%%%%%%%%%%%%%%%%%
%%  INICIO DO DOCUMENTO  %%
%%%%%%%%%%%%%%%%%%%%%%%%%%%
\begin{document}

% Retira espaço extra obsoleto entre as frases.
\frenchspacing

% ----------------------------------------------------------
% ELEMENTOS PRÉ-TEXTUAIS (Capa, Resumo, Abstract, etc.)
% ----------------------------------------------------------
\pretextual

% Capa
\include{pretextual/capa}



% Folha de rosto (o * indica que haverá a ficha bibliográfica)
%extras/folhaderosto é uma reimplementação do template abntex, que não precisa ser editado
\include{extras/folhaderosto} 
\imprimirfolhaderosto*


% Imprimir Ficha Catalografica
\include{pretextual/catalografica}

% Inserir Folha de Aprovação
\include{pretextual/assinaturas}

% Dedicatória
\include{pretextual/dedicatoria}

% Agradecimentos
\include{pretextual/agradecimentos}

% Epígrafe
\include{pretextual/epigrafe}

% Resumo e Abstract
\include{pretextual/resumos}

% Lista de ilustrações
\setlength{\afterchapskip}{\baselineskip}
\pdfbookmark[0]{\listfigurename}{lof}
\listoffigures*
\cleardoublepage


% Lista de tabelas
\setlength{\afterchapskip}{\baselineskip}
\pdfbookmark[0]{\listtablename}{lot}
\listoftables*
\cleardoublepage

% Lista de siglas - opcional podendo ser comentado/descomentado
\begin{siglas}
  \item[ABNT] Associação Brasileira de Normas Técnicas
  \item[abnTeX] ABsurdas Normas para TeX
\end{siglas}

% Lista de símbolos - opcional podendo ser comentado/descomentado
\begin{simbolos}
  \item[$ \Gamma $] Letra grega Gama
  \item[$ \Lambda $] Lambda
  \item[$ \zeta $] Letra grega minúscula zeta
  \item[$ \in $] Pertence
\end{simbolos}

% Inserir o SUMÁRIO
\pdfbookmark[0]{\contentsname}{toc}
\tableofcontents*
\cleardoublepage

% ----------------------------------------------------------
% ELEMENTOS TEXTUAIS (Capítulos)
% ----------------------------------------------------------
\textual

% Inclui cada capitulo da Monografia
\include{capitulos/introducao}
\include{capitulos/capitulo2}
\include{capitulos/capitulo3}
\include{capitulos/capitulo4}
\include{capitulos/conclusao}


% ----------------------------------------------------------
% ELEMENTOS PÓS-TEXTUAIS (Referências, Glossário, Apêndices)
% ----------------------------------------------------------
\postextual

% Referências bibliográficas
\citeoption{abnt-full-initials=yes}
\citeoption{abnt-emphasize=bf}
\citeoption{abnt-last-names=abnt}
\bibliography{bibliografia}
% Para montar seu arquivo "bibliografia.bib" você pode usar o http://bibdesk.sourceforge.net/ ou pode usar online o http://truben.no/latex/bibtex/


% Glossário (Consulte o manual)
% use apenas se tiver 100% de certeza do que está fazendo
%\glossary

% Apêndices
\include{postextual/apendices}

% Anexos
\include{postextual/anexos}

% Índice remissivo (Consultar manual)
% use apenas se tiver 100% de certeza do que está fazendo
%\phantompart
%\printindex

\end{document}
